\section{FOLIO Few-Shot Prompts}
\label{sec:appendix_folio_prompts}
The methodologies we investigate do not require any finetuning on domain-specific data.
Instead, we use in-context learning (ICL) with pretrained models.
We prompt the model with a set of instructions and 1-8 ICL examples, which adhere to a structured text format designed to scaffold generations and ease post-processing.
In particular, we begin each ICL example with each of the NL premises wrapped in an HTML-style tag \texttt{<PREMISES>\dots</PREMISES>} followed by the NL conclusion wrapped in \texttt{<CONCLUSION>\dots</CONCLUSION>}.
The requisite evaluation steps for each evaluation paradigm are then outlined in a subsequent section wrapped \texttt{<EVALUATE>\dots</EVALUATE>}.
Following the inclusion of ICL examples, a test example is added, with the \texttt{<PREMISES>} and \texttt{<CONCLUSION>} sections.
Then, the \texttt{<EVALUATE>} tag is opened, and the LM is allowed to proceed with causal generation until the \texttt{</EVALUATE>} tag is generated.
Upon generation of this stop token, the \texttt{<EVALUATE>} block is segmented for post-processing according to the method being evaluated (\texttt{\{na\"ive, scratchpad, chain-of-thought, neuro-symbolic\}}).

For the few-shot examples, we use samples from the publicly available FOLIO training set.
We select a set of diverse samples that are balanced across labels.
Since the FOLIO training set does not come with FOL expressions for the conclusions or chain of thought prompts, we manually add both for each sample.
For the $k$-shot setting ($k$<8), we use the first $k$ samples from the following list of sample indices: 126, 24, 61, 276, 149, 262, 264, 684.
Here, sample $i$ refers to the $i$th line in \url{https://github.com/Yale-LILY/FOLIO/blob/main/data/v0.0/folio-train.jsonl}.
We do not optimize for the choice of few-shot examples, and this is the only set of examples we evaluated with, so it is likely that there exist better choices for few-shot examples that would lead to improved performance across the board.
\subsection{FOLIO, 1-shot (baseline)}
\begin{lstlisting}[caption={todo},captionpos=b, breaklines=true]
The following is a first-order logic (FOL) problem.
The problem is to determine whether the conclusion follows from the premises.
The premises are given in the form of a set of first-order logic sentences.
The conclusion is given in the form of a single first-order logic sentence.
The task is to evaluate the conclusion as 'True', 'False', or 'Uncertain' given the premises.


<PREMISES>
All dispensable things are environment-friendly.
All woodware is dispensable.
All paper is woodware.
No good things are bad.
All environment-friendly things are good.
A worksheet is either paper or is environment-friendly.
</PREMISES>
<CONCLUSION>
A worksheet is not dispensable.
</CONCLUSION>
<EVALUATE>
Uncertain
</EVALUATE>

<PREMISES>
...premises for sample here, one premise per line
</PREMISES>
<CONCLUSION>
...conclusion for sample here
</CONCLUSION>
<EVALUATE>
\end{lstlisting}

\subsection{FOLIO, 1-shot (scratchpad)}

\begin{lstlisting}[breaklines=true]
The following is a first-order logic (FOL) problem.
The problem is to determine whether the conclusion follows from the premises.
The premises are given in the form of a set of first-order logic sentences.
The conclusion is given in the form of a single first-order logic sentence.
The task is to translate each of the premises and conclusions into FOL expressions, and then to evaluate the conclusion as 'True', 'False', or 'Uncertain' given the premises.


<PREMISES>
All dispensable things are environment-friendly.
All woodware is dispensable.
All paper is woodware.
No good things are bad.
All environment-friendly things are good.
A worksheet is either paper or is environment-friendly.
</PREMISES>
<CONCLUSION>
A worksheet is not dispensable.
</CONCLUSION>
<EVALUATE>
TEXT:	All dispensable things are environment-friendly.
FOL:	all x. (Dispensable(x) -> EnvironmentFriendly(x))
TEXT:	All woodware is dispensable.
FOL:	all x. (Woodware(x) -> Dispensable(x))
TEXT:	All paper is woodware.
FOL:	all x. (Paper(x) -> Woodware(x))
TEXT:	No good things are bad.
FOL:	all x. (Good(x) -> -Bad(x))
TEXT:	All environment-friendly things are good.
FOL:	all x. (EnvironmentFriendly(x) -> Good(x))
TEXT:	A worksheet is either paper or is environment-friendly.
FOL:	((Paper(Worksheet) & -EnvironmentFriendly(Worksheet)) | (-Paper(Worksheet) & EnvironmentFriendly(Worksheet)))
TEXT:	A worksheet is not dispensable.
FOL:	-Dispensable(Worksheet)
ANSWER:	Uncertain
</EVALUATE>

<PREMISES>
...premises for sample here, one premise per line
</PREMISES>
<CONCLUSION>
...conclusion for sample here
</CONCLUSION>
<EVALUATE>
\end{lstlisting}

\subsection{FOLIO, 1-shot (chain of thought)}

\begin{lstlisting}[breaklines=true]
The following is a first-order logic (FOL) problem.
The problem is to determine whether the conclusion follows from the premises.
The premises are given in the form of a set of first-order logic sentences.
The conclusion is given in the form of a single first-order logic sentence.
The task is to translate each of the premises and conclusions into FOL expressions, 


<PREMISES>
All dispensable things are environment-friendly.
All woodware is dispensable.
All paper is woodware.
No good things are bad.
All environment-friendly things are good.
A worksheet is either paper or is environment-friendly.
</PREMISES>
<CONCLUSION>
A worksheet is not dispensable.
</CONCLUSION>
<EVALUATE>
Let's think step by step. We want to evaluate if a worksheet is not dispensable. From premise 6, we know that a worksheet is either paper or is environment-friendly. If it is paper, then from premise 3, a worksheet is woodware, and from premise 2, a worksheet is dispensable. If it is environment-friendly, we know it is good from premise 5, but we know nothing about whether it is dispensable. Therefore, we don't know if a worksheet is dispensible or not, so the statement is uncertain.
ANSWER:	Uncertain
</EVALUATE>

<PREMISES>
...premises for sample here, one premise per line
</PREMISES>
<CONCLUSION>
...conclusion for sample here
</CONCLUSION>
<EVALUATE>
\end{lstlisting}

\subsection{FOLIO, 1-shot (neurosymbolic)} \label{sec:folio1shotexample}
\begin{lstlisting}[breaklines=true]
The following is a first-order logic (FOL) problem.
The problem is to determine whether the conclusion follows from the premises.
The premises are given in the form of a set of first-order logic sentences.
The conclusion is given in the form of a single first-order logic sentence.
The task is to translate each of the premises and conclusions into FOL expressions, so that the expressions can be evaluated by a theorem solver to determine whether the conclusion follows from the premises.
Expressions should be adhere to the format of the Python NLTK package logic module.


<PREMISES>
All dispensable things are environment-friendly.
All woodware is dispensable.
All paper is woodware.
No good things are bad.
All environment-friendly things are good.
A worksheet is either paper or is environment-friendly.
</PREMISES>
<CONCLUSION>
A worksheet is not dispensable.
</CONCLUSION>
<EVALUATE>
TEXT:	All dispensable things are environment-friendly.
FOL:	all x. (Dispensable(x) -> EnvironmentFriendly(x))
TEXT:	All woodware is dispensable.
FOL:	all x. (Woodware(x) -> Dispensable(x))
TEXT:	All paper is woodware.
FOL:	all x. (Paper(x) -> Woodware(x))
TEXT:	No good things are bad.
FOL:	all x. (Good(x) -> -Bad(x))
TEXT:	All environment-friendly things are good.
FOL:	all x. (EnvironmentFriendly(x) -> Good(x))
TEXT:	A worksheet is either paper or is environment-friendly.
FOL:	((Paper(Worksheet) & -EnvironmentFriendly(Worksheet)) | (-Paper(Worksheet) & EnvironmentFriendly(Worksheet)))
TEXT:	A worksheet is not dispensable.
FOL:	-Dispensable(Worksheet)
</EVALUATE>

<PREMISES>
...premises for sample here, one premise per line
</PREMISES>
<CONCLUSION>
...conclusion for sample here
</CONCLUSION>
<EVALUATE>
\end{lstlisting}